\lysh{14975}{4}
\begin{alist}
\item Η τριγωνομετρική συνάρτηση $f(x) = \rho \cdot \hm (\omega x)$, όπου $\rho, \omega > 0$ έχει μέγιστη τιμή $\rho$, ελάχιστη τιμή $-\rho$ και περίοδο $T = \dfrac{2\pi}{\omega}$.
Ως εκ τούτου,
Το $y_0 = (OA) = (OB) = \rho = 0{,}2$ μέτρα.
Η συνάρτηση $y(t)$ έχει μέγιστη τιμή $1 + \rho = 1{,}2$, ελάχιστη τιμή $1 - \rho = 0{,}8$ και η απόσταση μεταξύ των δύο ακραίων θέσεων Α και Β της ταλάντωσης είναι:
$|1{,}2 - 0{,}8| = 0{,}4$ μέτρα.
\item η περίοδος της συνάρτησης $y(t)$ είναι: $T = \dfrac{2\pi}{\frac{\pi}{2}} \Leftrightarrow T = 4$.
\item Ο πίνακας τιμών για τη συνάρτηση $y(t) = 1 + 0{,}2 \cdot \hm \left(\dfrac{\pi}{2}t\right)$ για $t \in [0, 4]$, είναι:
\begin{center}
\begin{myhtblr}{}
$t$ & $0$ & $1$ & $2$ & $3$ & $4$ \\
$y(t)$ & $1$ & $1{,}2$ & $1$ & $0{,}8$ & $1$ \\
\end{myhtblr}
\end{center}
Είναι $y_{max} = 1{,}2$ και $y_{min} = 0{,}8$ και η γραφική παράσταση της συνάρτησης:
\begin{center}
\begin{tikzpicture}[x=1cm,y=1cm,scale=3]
\draw[xstep=0.5cm,ystep=0.2cm,gray!30,very thin] (-0.7,-0.3) grid (4.5,1.9);
\draw[->] (-0.5,0) -- (4.5,0) node[above] {$t$};
\draw[->] (0,-0.2) -- (0,1.9) node[right] {$y$};
\foreach \x/\xtext in {0.5/0{,}5, 1/1, 1.5/1{,}5, 2/2, 2.5/2{,}5, 3/3, 3.5/3{,}5, 4/4} \draw (\x cm,1pt) -- (\x cm,-1pt) node[anchor=north] {\scriptsize \xtext};
\foreach \y/\ytext in {0.2/0{,}2, 0.4/0{,}4, 0.6/0{,}6, 0.8/0{,}8, 1/1, 1.2/1{,}2, 1.4/1{,}4, 1.6/1{,}6, 1.8/1{,}8} \draw (1pt,\y cm) -- (-1pt,\y cm) node[anchor=east] {\scriptsize \ytext};
\node[below left] at (0,0) {\scriptsize $0$};
\node[left] at (0,-0.2) {\scriptsize $-0{,}2$};
\node[below] at (-0.5,0) {\scriptsize $-0{,}5$};
\draw[maincolor!60!black, thick, domain=0:4, samples=100] plot (\x,{1 + 0.2*sin(90*\x)});
\node[maincolor!60!black] at (2, 1.5) {$y(t) = 1 + 0{,}2 \, \hm  \dfrac{\pi}{2} t$};
\filldraw[maincolor!60!black] (0,1) circle (.5pt);
\filldraw[maincolor!60!black] (1,1.2) circle (.5pt);
\filldraw[maincolor!60!black] (2,1) circle (.5pt);
\filldraw[maincolor!60!black] (3,0.8) circle (.5pt);
\filldraw[maincolor!60!black] (4,1) circle (.5pt);
\end{tikzpicture}
\end{center}
\item Ζητάμε ουσιαστικά να βρούμε ποια χρονική στιγμή $t \in [0, 2]$, είναι $y(t) = 1{,}1$.
\end{alist}

